% Structure: (1) problem, (2) method, (3) headline result, (4) mechanism.
% The headline is the controlled margin; the anchor finding SUPPORTS it and
% must not be allowed to become the headline again.
Behaviour-cloned robot policies are capped by their demonstrations, and the
usual way past that cap is to train a critic and fine-tune against it. This
reliably goes wrong: the policy drifts away from the data the critic was fit
to, its score rises, and its actual success collapses.

We present an actor update that does not. A frozen base proposes a cloud of
candidate actions; each is displaced along the critic's gradient; a restoring
force returns any candidate that has strayed beyond a radius $\rho$ of the
base; and a small residual head is regressed onto the displaced candidates.
Nothing else in the pipeline changes --- the critic, replay, demonstrations,
budget and evaluation are those of the state-of-the-art residual fine-tuner,
byte for byte, so the only variable is the actor update.

Under that control, the update gains $+9.4$ points on the eight hardest
LIBERO-90 tasks where backpropagating $-Q$ gains $+0.9$, and lifts robomimic
square from $0.382$ to $0.91$--$0.93$. It reaches the published
flow-matching pipeline's performance from a base less than half as strong,
and transfers unchanged to that flow-matching base.

We then isolate which part is responsible, by switching each guard off
independently. The restoring force is: without it the policy ends below its
untuned base. A bound on step size, the more obvious remedy, is not --- on
its own it leaves the policy at zero success, because a drift made of many
individually permitted steps is still a drift. The reason appears to be that
a restoring force with a dead zone is silent exactly where the critic is
reliable, whereas the standard always-on behaviour-cloning penalty opposes
the critic everywhere and must trade containment against progress with one
coefficient.
